\label{sec:literature}
% ---- styles for the literature-study figures ----
\tikzset{
  litbox/.style={draw,rounded corners=2pt,fill=gray!12,align=center,font=\footnotesize,minimum height=7mm,inner sep=3pt},
  litio/.style={draw,rounded corners=6pt,fill=white,align=center,font=\footnotesize,minimum height=7mm,inner sep=3pt},
  litarr/.style={-{Stealth[length=2.0mm]},thick},
  litlab/.style={font=\scriptsize,inner sep=1.5pt,align=center}
}

\new{Text in red is not the author's own wording: it marks the pen-mapping description, the intended contribution, and the individual sentences that were corrected for factual accuracy. All remaining text is the author's own.}

\subsection{Existing concepts}
\label{sec:lit-background}

\subsubsection{A need exists to be able to decode the writer's style and written text from an image}
\label{sec:lit-hwr}

The ability of a computer or system to identify written text after being trained on large amounts of data is known as HWR (Handwriting Recognition) or HTR (Handwritten Text Recognition)~\cite[p.~39]{chen2023}.
One type of handwritten text recognition is offline recognition, which means the data being identified by the system is static and captured post-completion of the page of handwritten text. This is ideal for identifying the text from a fixed image of handwritten content by utilizing image processing techniques and algorithms to identify the text that was written~\cite{ma2023}. The other type is online recognition, which is a more dynamic process that identifies text during the creation of the handwriting, adding the variable of time to the classification of text. The addition of a time variable includes data about the order of strokes to create text as well as pen-up or pen-down movements, which can influence the type and ability of reproduction and handwriting synthesis methods.

Every stage of the process to achieve accurate handwritten text recognition requires a deep understanding of the different techniques available:

\textbf{Datasets and resolution:} Multiple datasets have typically been used for training handwritten text recognition modules, for example, the EMNIST, MNIST and IAM datasets~\cite[p.~2]{mamun2024}. The EMNIST (Extended Modified National Institute of Standards and Technology) dataset consists of approximately 814\,255 handwritten characters and numbers, each being 28 by 28 pixels in size~\cite[p.~40]{chen2023}. There is only a single, disconnected character per image, which is used in many existing HTR models, but these models cannot effectively decipher connected handwritten words, like cursive words. The IAM handwriting dataset provides much more variety, as it consists of 5\,685 complete sentences, with a total of 115\,320 words, gathered from 657 different authors, consisting of 26 letters and 27 punctuation marks~\cite{iamdb}; the paper that introduced the database reports an earlier state of it, 1\,066 forms by about 400 writers~\cite[p.~39]{marti2002}. The images from the IAM dataset were scanned at 300 dpi with 8 bits per pixel in grayscale~\cite{iamdb}. The dataset used influences the accuracy for the classification of text with higher variation, as well as defines the resolution of the camera required to capture the handwritten text to classify. \new{The IAM distribution of cropped lines is not organised by author, so the writer is not evident from the line images alone, while the distribution of full pages is arranged so that each page belongs to a known writer; the database does record the writer of every form in its metadata, so the two can be joined~\cite{marti2002}.}

\textbf{Image-conditioning:} The image-conditioning of handwriting pages includes multiple techniques to be performed on the handwritten text before being passed into the classification and feature recognition algorithms. This is used to account for dimensional adjustments, removal of inconsistencies, and normalization of the data, to make the images more appropriate for text separation and classification~\cite{hemanth2022}. Normalization is used to increase the accuracy of HTR models by reducing the effect that unwanted features of the data could have on the performance of the model. There are multiple different normalization techniques such as min-max scaling, z-score normalization, and log-scaling. Encoding techniques are also used in HTR model pre-processing to get rid of unintentional biases~\cite[pp.~40--41]{chen2023}. Binarization is used to differentiate the background page and obscurities from the actual ink of the written text, computed by pixel. One method of binarization is done from a global ink threshold taken from the grey-level histogram, which does not require any parameters like a localized shadow level radius, and is used for ideal, evenly lit page scans. The other method involves a localized threshold computed per pixel out of the statistics of a window around it, which costs more in terms of computation time and power but can tolerate a lighting gradient or shadow. Deskewing is a technique used to fix the rotation of scanned pages, by rotating the page through certain angle steps until the horizontal projection of the ink is most sharply peaked, which occurs when the lines of text all mostly lie horizontally. Deslanting is used to fix the lean of the individual characters \new{by applying a horizontal shear that shifts each row sideways in proportion to its height}, forcing more pixels to be aligned within the same column rather than slanted over multiple columns.

\textbf{Text Separation:} Text separation is the process of breaking up a page of handwritten text into smaller segments or blocks that can be used in feature extraction and neural network classifiers. There are three methods that existing models employ to separate text: line separation, individual character separation or word separation. All three methods involve breaking down paragraphs of handwritten text into smaller segments. The accuracy of the extraction of smaller segments heavily influences the effectiveness of the classifier. Some existing techniques for this are classic dissection techniques, recognition-based segmentation, hysteresis thresholding, and bounding box techniques~\cite[p.~41]{chen2023}. Line separation can be done by converting images to grayscale, obtaining inverted binary images, and then drawing bounding boxes over each line~\cite{hemanth2022}. Firstly, when lines are mostly straight and separated by a blank row, line separation can be performed by reading the minimums of the row ink profile (the sum of ink in each row) to form boundaries. When descenders of lines begin to overlap or merge, the straight row boundary separator is replaced with a boundary that bends around the descenders by performing a least-cost path traced between the ink lines, and around the descenders.
To achieve either word or character separation, the page must still be broken up into individual lines and then into separate words. Character separation just goes one step further and splits the word into individual characters.

\textbf{Types of Neural Networks:} Many different types of neural network models have been implemented in previous papers, including ANNs (artificial neural networks), CNNs (convolutional neural networks), RNNs (recurrent neural networks), and various combinations of them to achieve the best accuracy.

ANNs are created from the standard neural network architecture consisting of multiple neurons, weights, hidden layers, biases, etc., and can accurately perform pattern recognition for handwritten text. CNNs are built on the architecture of ANNs but also facilitate feature extraction and identifying smaller, more important information in text. This leads to a reduction of necessary learnable parameters and streamlines model training~\cite{ma2023}. RNNs contain the same fundamental structure as CNNs but also have the capability to consider sequential and dependent data. The ability to capture time-dependent patterns is useful for deciphering various handwritten fonts, and because RNNs can consider dependencies in the data, this can improve the accuracy of handwritten text recognition when considering common linguistic dependencies in language~\cite{ma2023}.
A combination of both CNNs and RNNs can significantly enhance feature extraction by using the CNN's pattern recognition and ability to identify diverse handwriting styles, and the RNN's ability to use textual dependencies for improved contextual understanding~\cite{kumar2024}. This combination for improved context awareness allows an entire line of text to be read without word or character separation. The convolutional layers are used to convert a cropped line of text into different feature columns, and the bidirectional recurrent layer reads the sequence of features in each direction, providing context of what comes before and after each character or feature. This implementation is used for unconstrained offline handwriting~\cite[p.~855]{graves2009}, \new{meaning natural handwriting that is not restricted to isolated characters or to a fixed vocabulary, and the full line crop must exist before classification can begin because the recurrent layer reads it from both ends}. In an RNN, the gradient is multiplied by the same weight matrix at every time step as it propagates backwards, so over many steps it shrinks and loses its effect (this is the vanishing gradient problem). Forget gates are added to bypass the vanishing effect and allow gradient errors to still retain an effect after many time steps. Putting convolutional layers before the recurrent stack allows the feature extraction and sequence reading to be trained as one network~\cite{shi2017}. The columns produced from the convolutional layers and max pooling tend to outnumber the number of characters to be classified, therefore \new{the CTC (connectionist temporal classification) loss~\cite[p.~370]{graves2006} is used, which adds a blank symbol and sums the probability of every column-to-character alignment that collapses to the transcription, the collapse being to squash repeated symbols first and then remove the blanks}. The arrangement is shown in Figure~\ref{fig:lit-line}.

\begin{figure}[!htbp]
\centering
\begin{tikzpicture}[x=1cm,y=1cm,font=\footnotesize,node distance=4.5mm]
\node[litio,text width=1.5cm] (img) {line image};
\node[litbox,right=of img,text width=1.7cm] (cnn) {convolutional\\layers};
\node[litbox,right=of cnn,text width=1.6cm] (cols) {$T$ feature\\columns};
\node[litbox,right=of cols,text width=1.9cm] (rnn) {recurrent layer,\\both directions};
\node[litbox,right=of rnn,text width=1.6cm] (pc) {symbol per\\column};
\node[litio,right=of pc,text width=1.1cm] (out) {\texttt{the}};
\foreach \a/\b in {img/cnn,cnn/cols,cols/rnn,rnn/pc,pc/out}{\draw[litarr] (\a) -- (\b);}
%% the two reading directions
\draw[litarr] ([yshift=3.2mm]rnn.north west) -- ([yshift=3.2mm]rnn.north east);
\draw[litarr] ([yshift=6.0mm]rnn.north east) -- ([yshift=6.0mm]rnn.north west);
%% the blank-augmented output and the collapse
\node[litlab,below=2.0mm of pc,font=\ttfamily\scriptsize] (sym) {-tt-h-ee};
\node[litlab,below=1.6mm of sym] {collapse repeats,\\drop the blanks};
\node[litlab,below=2.0mm of cols,text width=2.5cm] {$T$ far exceeds the\\number of characters,\\so the alignment is\\not known};
\end{tikzpicture}
\caption{Reading a whole text line without cutting out characters, after the arrangement of Kizilirmak and Yanikoglu~\cite[pp.~4--5]{kizilirmak2023}. The convolutional layers produce $T$ feature columns, a recurrent layer reads them in both directions, and a symbol including a blank is emitted per column; collapsing repeats and removing the blanks gives the transcription.}
\label{fig:lit-line}
\end{figure}

\textbf{Existing Models:} Two baseline models have been primarily considered in existing research papers: the LeNet-5 CNN model and the ResNet-34 CNN model. LeNet-5 is the most common of the two and was created by Yann LeCun in \new{1998}~\cite[Fig.~2]{lecun1998}. It is a CNN structure that consists of \new{seven} layers, including 3 convolutional layers, 2 pooling layers, and lastly, 2 fully connected layers. A fully connected layer refers to a portion of a typical ANN model, while the convolutional layers and pooling layers will be discussed later in this literature study. The input size to the model is 32 by 32 pixels, and the model uses a hyperbolic tangent activation function in the fully connected layers and a softmax activation function in the output layer. This model is popular for use in object detection and image classification~\cite{chen2023}.
ResNet-34 is also a popular CNN model known for its residual architecture. \new{The residual architecture adds skip connections that let the gradient bypass layers, which is what allows the network to be trained at that depth; separately, its weights from training on a large image set are commonly reused by transfer learning, which forms the foundation of the feature extraction in this model and allows for fast and efficient training from limited datasets}~\cite{ma2023}.

\textbf{Feature extraction:} The convolutional layers of a CNN are used to perform feature extraction, which means they can capture subtle patterns in different areas of a segment of text by sliding a convolutional kernel over the character or word image. Translation invariance is created by sliding the convolution kernel, which allows consistent patterns to be extracted from images~\cite{ma2023}. An $N$ by $N$ input layer is convolved with an $f$ by $f$ kernel and returns an $(N-f+1)$ by $(N-f+1)$ output. Feature maps are created by applying convolutions to different sections of the image being classified. By identifying common features, the kernel-sliding convolutional technique can produce more accurate classification. Each extracted feature is multiplied with a corresponding kernel value to produce a weighting, and then a bias and activation function are applied until an output feature map is obtained~\cite{yahya2021}.

An extended version of the standard CNN is the deep CNN, which refers to a CNN that can extract higher-complexity features like handwritten strokes. The problem is that a larger, deeper CNN creates a `vanishing gradient' problem where there are so many layers that during backpropagation, some layers never get their weightings changed. Thus, there is a limit to the number of layers a standard CNN can have, where adding more layers yields no further benefit. This is why a residual network is used to implement a deep CNN by using residual blocks and skip connections across various layers so that all layers will still update their weightings during backpropagation regardless of the number of layers~\cite{ma2023}.

After a singular round of feature selection, many parameters are still required to classify text with these high-level features. Max-pooling is used to downsample these feature maps until only the essential features remain, minimizing the number of parameters needed to classify handwritten characters and preventing overfitting~\cite{ji2021}. Figure~\ref{fig:lit-conv} shows one convolutional layer followed by one max-pooling layer, which is the unit the models are stacked from.

\begin{figure}[!htbp]
\centering
\begin{tikzpicture}[x=1cm,y=1cm,font=\footnotesize]
%% input N x N with the f x f kernel window
\draw[step=0.26,gray!50,thin] (0,0) grid (1.56,1.56);
\draw[thick] (0,0) rectangle (1.56,1.56);
\fill[gray!40,opacity=0.75] (0.26,1.04) rectangle (1.04,1.56);
\draw[very thick] (0.26,1.04) rectangle (1.04,1.56);
\node[litlab,below] at (0.78,-0.10) {input, $N\times N$};
\node[litlab,above] at (0.78,1.62) {$f\times f$ kernel slides};
%% arrow 1
\draw[litarr] (1.75,0.78) -- (2.60,0.78);
\node[litlab,above] at (2.17,0.86) {multiply,};
\node[litlab,below] at (2.17,0.72) {add bias};
%% feature map (N-f+1)
\draw[step=0.26,gray!50,thin] (2.80,0.13) grid (3.84,1.43);
\draw[thick] (2.80,0.13) rectangle (3.84,1.43);
\fill[gray!40,opacity=0.75] (2.80,1.17) rectangle (3.06,1.43);
\node[litlab,below] at (3.32,0.03) {feature map,};
\node[litlab,below] at (3.32,-0.22) {$(N-f+1)\times(N-f+1)$};
\node[litlab,above] at (3.32,1.49) {activation applied};
%% arrow 2
\draw[litarr] (4.03,0.78) -- (4.88,0.78);
\node[litlab,above] at (4.45,0.86) {max-pool};
\node[litlab,below] at (4.45,0.72) {$2\times2$};
%% pooled map
\draw[step=0.26,gray!50,thin] (5.08,0.39) grid (5.60,1.17);
\draw[thick] (5.08,0.39) rectangle (5.60,1.17);
\node[litlab,below] at (5.34,0.29) {downsampled};
\node[litlab,below] at (5.34,0.04) {feature map};
%% stacking note
\draw[litarr,dashed] (5.79,0.78) -- (6.55,0.78);
\node[litlab,right,align=left] at (6.60,0.78) {repeated, each stage\\ holding more feature\\ maps than the last};
\end{tikzpicture}
\caption{One convolutional layer and one max-pooling layer. The kernel slides over every position of the input, so an $N\times N$ input gives an $(N-f+1)\times(N-f+1)$ feature map, which pooling then downsamples.}
\label{fig:lit-conv}
\end{figure}

\textbf{Training:} The way networks are trained corresponds to the architecture they are built from. The weights are initially randomized with a variance determined by the size of each layer so the signal remains constant as it passes through more layers~\cite{he2015}. Batch normalization is the process of taking the outputs of the current batch of training data and subtracting the average and dividing by their \new{standard deviation}~\cite[Alg.~1]{ioffe2015}. Having more sensibly ranged values allows a higher learning rate to be used safely. Dropout is a technique used to reduce overfitting which removes a random subset of \new{a layer's outputs on each training step}~\cite[p.~1930]{srivastava2014}. This introduces variation and prevents the network from depending on any single feature. Another method to increase variation is dataset enlargement which includes applying shears, small rotations, slants and random smooth warps etc. to simulate more variation in training text~\cite{simard2003}. Learning rate reduction is a technique that is used to adjust the learning rate over time to allow for large changes initially to speed up training, and fine learning later on to prevent overshooting. Final accuracy can be calculated as the character error rate (which refers to the number of mismatches in characters compared to the ideal result) or the word error rate, which is calculated from word mismatches. This is usually significantly higher than the character error rate as one character mismatch ruins the whole word~\cite[p.~1]{kizilirmak2023}.

\textbf{Techniques for writer identification:} There are multiple approaches to identifying the author of a written piece of text. The first way is through feature extraction, which measures text-independent features like slant, height, stroke width, word and character spacing, and connectivity between characters. The author is determined by finding which writer possesses the feature metrics closest to that of the text measured. Another feature that can be extracted is the \new{distribution of edge directions in the ink, which gives the} `texture' of the written sentence~\cite{bulacu2007}. The other method of writer identification is using a neural network trained off the author's handwriting, which can use a feature extractor used for text recognition~\cite{ma2023} instead of the hardcoded methods and metrics for quantifying features. Much more content and pages are required to train the author classifiers in this case compared to creating specific feature extracting methods.

\subsubsection{A need exists to reproduce text on paper accurately and fast in any writer's style}
\label{sec:lit-synth}

There are various existing text reproduction methods and linear actuation techniques which can influence the accuracy and speed of text reproduction.

\textbf{Handwriting synthesis in a writer's style:} Before being able to reproduce handwriting through hardware, the desired text to be reproduced has to still be created in the particular author's style. \new{Four} methods will be discussed on how to achieve this:

\emph{Recurrent stroke level synthesis:} A recurrent neural network is used to predict not the characters to be written, but rather the pen movement that must be taken next, with small offset predictions at a time~\cite[p.~1]{graves2013}. It creates a spread of possible answers as no letter is written identically. The network is trained off an online recorded sample of someone's handwriting. Since it predicts pen movements the outputs are already presented as movements which a plotter requires. Training a network to produce accurate outputs requires hundreds of lines of online data.

\emph{Adversarial image synthesis:} This method creates an image of text using a GAN (generative adversarial network), which trains a text generator and discriminator to identify fake generated text from real handwritten text~\cite[p.~273]{kang2020}. Both the generator and discriminator improve per round until the generator can produce well-reproduced text and the discriminator can no longer tell the difference. It is called `adversarial' as two networks are trained simultaneously, one to generate text and the other to judge the text. This requires a significant amount of time and compute to train, but the final system produces high quality recognisable text \new{in a single pass once trained}. A text and writer classifier are still required to verify the generated text along the way as well. This method produces a raster image which must be converted to pen paths for a plotter to take.

\emph{Diffusion-based synthesis:} This method generates a fake text raster image by reversing a gradual noise-generating process~\cite[Sec.~1]{dai2024}. The system trains by slowly adding noise to handwritten text, remembering the way noise was added, until the final image is pure random noise. Then to generate a new line of text, the system starts with complete random noise and denoises it until an accurate line of text is reproduced. \new{Only one reference image of the target writer is needed at generation time, although training the model still requires a large set of labelled handwriting images}, yet it requires many steps and time to denoise the random noise image. This method also produces a raster image which must still be converted to pen stroke mappings.

\emph{Glyph-library based synthesis:} This method does not use any generators to create text, but rather assembles words from shapes or glyphs taken out of existing handwritten pages. Haines \emph{et al.}~\cite[p.~26:1]{haines2016} describe a technique that stores a library of labelled glyphs per each user, and then creates the specified sentence by choosing the specific glyphs and then assembling them according to which variant of the glyph the writer usually used in the requested position, and the level of connectedness of the text, the texture of the line of ink, and the spacing of the writing. It requires a small number of handwritten pages per author to create decent glyph libraries, and the main complexities of the method only involve connecting one glyph to another and selecting the most appropriate glyphs.

\textbf{Assessing generated handwritten text:} Accuracy of generated handwriting can be measured through performing image-level distances and correlation to real handwritten text, for example an IoU test (Intersection over Union) which measures the overlap between a predicted region and the true region. Another method is by passing it through an independent text and author classifier and reporting the accuracy of the reproduced text. A survey can also be conducted for human validation of what text appears real and what appears reproduced.

\textbf{Method of reproduction:} There are three main methods to reproduce handwritten text on paper: Cartesian systems, serial-chain robots, and parallel-chain robots. A Cartesian system consists of a gantry with x\,--\,y movement. It has restricted movement in only two directions using linear actuation. Serial-chain robots are the most common in industrial settings for automated movement. They consist of multiple linkages chained together, one component connected to the next in a single chain, like a typical robot arm. The third method is parallel-chain robots, which consist of multiple serial-chain linkages that can precisely control automated movement~\cite{teikari2016}.

\textbf{Pen mapping techniques:} \new{A pen path is stored as a list of points, and a drawn stroke usually has far more points than a plotter needs. Pen mapping is the step that reduces this list while keeping the shape of the stroke. The Douglas-Peucker algorithm is the common choice~\cite{teikari2016}: it draws a straight line between the two ends of a stroke, finds the point that sits furthest from that line, and keeps it only if it is further away than a chosen tolerance, then repeats on each half~\cite{douglas1973}. A small tolerance keeps many points and follows the stroke closely, while a large tolerance keeps few points and cuts corners, so the tolerance sets the trade-off between the file size and the accuracy of the drawn shape. The alternatives are fitting a smooth curve such as a spline through the points, which needs a fit to be solved, and resampling the stroke at a fixed spacing, which spaces the points evenly but ignores where the curves are. When the handwriting only exists as an image, the path has to be recovered first by thinning the ink down to a one-pixel centre line and then joining those pixels into strokes, which means guessing the drawing order wherever two strokes touch or cross.}

\textbf{Linear actuation types:} There are multiple methods of linear actuation that could be considered for a Cartesian plotting system. Some popular examples of linear actuation methods are electric, pneumatic, belt-drive, rack-and-pinion, and ball-screw systems. Electric and pneumatic methods are often used for industrial cases, but belt-drive, rack-and-pinion, and ball-screw linear actuators have been used in multiple handwriting reproduction systems due to their durability, reliability, and cost-effectiveness~\cite{iqs2024}.

\subsection{Intended contribution and discussion}
\label{sec:lit-application}
\label{sec:lit-embedded}

\new{The literature above covers each stage of the problem on its own, but it does not cover them as one machine. Recognition is reported on evenly lit scans or on isolated characters, writer identification is reported over sets of hundreds of writers rather than a small closed group, and handwriting synthesis is judged as an image rather than as something that has to be drawn with one pen. Almost all of this work is also built on libraries such as OpenCV~\cite[p.~41]{chen2023}, which hide what each stage actually costs in operations and memory. That cost matters here, because every stage of this project must run on a single-board computer.}

\new{The main contribution of this project is therefore the full integration of the system, from a photographed page to ink on paper, with every stage implemented from first principles and no machine-learning or image-processing libraries used anywhere in the pipeline. This covers the creation of the dataset, the conditioning and line separation of photographed pages, a line-level text recogniser, the identification of the author, the synthesis of new text in that author's style, and the generation of the pen paths and G-code that drive the gantry. Writing each stage by hand means the forward and backward passes of every layer, the loss functions and the optimisers are known quantities whose cost can be measured and reduced to fit the available hardware.}

\new{Two things set the approach apart from the work reviewed. The first is that the system classifies both the text that was written and the author who wrote it, using the same feature extractor for both, so that one page yields a transcription and an identity. The second is that the synthesised handwriting is produced as geometry rather than as an image, so it can be drawn directly without having to recover a pen path from pixels, and the recogniser and the author classifier are then reused to judge the synthesised output before it is committed to paper. The reasoning behind each of these choices, and the alternatives that were rejected, is given in Section~\ref{sec:design}.}
